\chapter{Capture et stockage des données tick}\label{ch:pl:capturing-and-storing-tick-data}

Le flux consolidé des options américaines doit transporter des centaines de milliards de messages un jour chargé (l'encadré daté de
la \cref{sec:pl:capturing-and-storing-tick-data:economics} donne le chiffre). Une société qui conserve chacun de ces messages pendant
cinq ans exploite en parallèle une activité de stockage, et la première décision qu'elle prend n'est pas quelle base de données acheter mais
ce qu'il faut exactement conserver~: les paquets tels qu'ils sont arrivés, avec leurs doublons, leurs trous et leurs heures d'arrivée,
ou les événements qu'ils désignaient, nettoyés, arbitrés et disposés pour la lecture. Ce chapitre capture dix minutes chargées de la
bourse simulée du livre~10 sur ses deux lignes de flux, les transforme en enregistrements d'événements de la société, mesure le coût de chaque disposition en octets, et met à l'échelle le
résultat, avec des volumes de messages et des prix cités, à la question du titre de son problème du week-end~: un pétaoctet par an~?

\section{Qu'est-ce que les données tick et où sont-elles capturées}

\begin{definition}[Données tick]\label{def:pl:capturing-and-storing-tick-data:tick}
Les \emph{données tick}\index{données tick} sont l'enregistrement de chaque événement qu'une plateforme publie sur ses instruments --- ajouts
d'ordres, modifications, annulations et exécutions, transactions, cotations, messages d'état et de mesures de négociation --- avec les horodatages que
la plateforme leur a apposés et l'heure à laquelle la société les a reçus.
\end{definition}

Les \omterm{def:pl:capturing-and-storing-tick-data:tick}{données tick} sont l'enregistrement le plus fin d'un marché que puisse conserver une société, et tout ce qui est plus grossier en
découle~: les barres (livre~7, chapitre~2), les variables de carnet d'ordres (chapitre~8 de ce livre), les rejeux de son chapitre~18. Ce qui l'est moins, c'est que «~les~»
\omterm{def:pl:capturing-and-storing-tick-data:tick}{données tick} d'une journée dépendent de l'endroit où elles ont été capturées. La \Cref{fig:pl:capturing-and-storing-tick-data:points} montre les trois
points de capture usuels sur le parcours du traitement des données de marché du livre~13.

\begin{omfigure}
\begin{tikzpicture}[font=\scriptsize,>=Stealth,
  box/.style={draw,rounded corners=2pt,fill=omDef!10,minimum height=0.7cm,inner sep=3pt,align=center},
  cpt/.style={draw,rounded corners=2pt,fill=omThm!12,minimum height=0.6cm,inner sep=3pt,align=center}]
\node[box] (a) at (0,0.6) {ligne A};
\node[box] (b) at (0,-0.6) {ligne B};
\node[box] (fh) at (3.2,0) {feed handler\\arbitrage, trous};
\node[box] (bk) at (7.0,0) {constructeur de\\carnet, stratégie};
\node[box] (st) at (10.9,0) {vue agrégée\\(écrans, risque)};
\draw[->] (a) -- (fh);
\draw[->] (b) -- (fh);
\draw[->] (fh) -- node[above,font=\tiny]{événements} (bk);
\draw[->] (bk) -- node[above,font=\tiny]{mises à jour} (st);
\node[cpt] (c1) at (1.4,-1.9) {1. au fil~:\\paquets, les deux lignes};
\node[cpt] (c2) at (4.9,-1.9) {2. après arbitrage~:\\un événement par message};
\node[cpt] (c3) at (9.0,-1.9) {3. après conflation~:\\ce sur quoi on a agi};
\draw[->,dashed,omThm] (1.4,0) -- (c1);
\draw[->,dashed,omThm] (4.9,0) -- (c2);
\draw[->,dashed,omThm] (9.0,0) -- (c3);
\end{tikzpicture}

\omcaption{Trois points de capture sur le parcours des données de marché. Au fil (une prise réseau ou un boîtier de capture, livre~14,
chapitre~5), les paquets des deux lignes sont conservés tels qu'ils sont arrivés~; après l'arbitrage du feed handler, chaque message n'apparaît qu'une fois,
dans l'ordre de la séquence~; après la conflation, seules restent les mises à jour qu'un consommateur a effectivement vues.}\label{fig:pl:capturing-and-storing-tick-data:points}
\end{omfigure}

Chaque point répond à des questions différentes. Seule la capture au fil peut dire quelle ligne a perdu quel paquet et à quel moment chaque paquet
est arrivé, ce dont a besoin une enquête sur la latence ou sur le réseau. Seule la capture après arbitrage restitue le marché tel que le carnet de la société
l'a vu, message par message. Seule la capture après conflation correspond à ce sur quoi un consommateur lent (livre~13, chapitre~12)
a effectivement agi. La recherche veut la deuxième~; les opérations et la conformité ont parfois besoin de la première ou de la troisième.

\begin{definition}[Capture brute]\label{def:pl:capturing-and-storing-tick-data:raw}
Une \emph{capture brute}\index{capture brute} est la suite des octets reçus en un point de capture, conservée sans interprétation,
chaque unité (un paquet, une trame) étant accompagnée de l'heure à laquelle la société l'a reçue.
\end{definition}

La capture du chapitre se fait au premier point, dans le format de fichier enregistré du simulateur du livre~10~: pour chaque paquet
reçu sur une ligne, son heure de réception (8~octets), sa longueur (4~octets) et le paquet MoldUDP64 lui-même. Dix minutes de la phase d'ouverture chargée du simulateur,
deux instruments, avec ses perturbations simulées --- pertes indépendantes sur chaque ligne, rafales de pertes, paquets dupliqués, gigue (jitter) et
une panne de cinq secondes sur la ligne B --- donnent 15.8~Mo sur la ligne A et 15.6~Mo sur la ligne B.

\section{Capture brute contre stockage normalisé}

Pour reconstituer à partir de la capture le marché que le livre a vu, les deux lignes sont fusionnées par numéro de séquence~: la première copie de chaque message
l'emporte, les copies suivantes sont des doublons, et un trou reste ouvert jusqu'à ce que l'autre ligne ait eu le temps de le combler. Ce que ni
l'une ni l'autre ligne ne comble constitue un trou dans l'enregistrement~; un feed handler en temps réel demanderait le service de retransmission de la plateforme ou reconstruirait à partir d'un instantané (snapshot)
(livre~13, chapitre~18), tandis qu'une capture destinée à la recherche le consigne et passe à la suite.
La \Cref{lst:pl:capturing-and-storing-tick-data:arb} est l'arbitrage de \texttt{firm.tickcap}.

\omcode{code/firm/tickcap/firm_tickcap.py}{53}{87}{Arbitrage de lignes pour le stockage~: la première copie l'emporte, les doublons sont comptés, et un
trou que ni l'une ni l'autre ligne ne comble dans le délai imparti est consigné puis ignoré.}\label{lst:pl:capturing-and-storing-tick-data:arb}

Sur la session de dix minutes, les deux lignes transportent 503\,024 paquets. L'arbitrage délivre 258\,535 messages, écarte 255\,899
doublons (presque chaque paquet arrive deux fois, une fois par ligne) et consigne quatre trous, soit huit messages au total, perdus sur les deux lignes au
même instant. Chaque message délivré est ensuite normalisé en enregistrement d'événement de la société~: l'heure de réception suivie de l'événement de 48~octets du
feed handler du livre~13 (nature, côté, instrument, numéro de séquence, heure d'échange, références d'ordres, prix, quantité), soit 56~octets au
total. Le stockage et le chemin de négociation partagent donc une seule disposition, et un test vérifie que les événements de la capture sont, message par
message, ceux du feed handler.

\begin{proposition}[La brute suffit, la normalisée non]\label{prop:pl:capturing-and-storing-tick-data:sufficient}
Les enregistrements normalisés d'une session sont une fonction de sa \omterm{def:pl:capturing-and-storing-tick-data:raw}{capture brute} (et du code du normaliseur). La \omterm{def:pl:capturing-and-storing-tick-data:raw}{capture brute}
n'est pas une fonction des enregistrements normalisés.
\end{proposition}

\begin{proof}
La première partie relève de la construction~: l'arbitrage et la normalisation sont des fonctions déterministes des paquets captés et
de leurs heures de réception. Pour la seconde, deux captures qui ne diffèrent que par ce que la normalisation écarte --- la ligne qui a livré un message
en premier, les doublons, le découpage en paquets, l'heure de réception d'une copie ultérieure, les octets d'un type de message que le normaliseur ne connaît pas
--- donnent les mêmes enregistrements.
\end{proof}

La proposition est l'argument entier en faveur de la conservation des \omterm{def:pl:capturing-and-storing-tick-data:raw}{captures brutes}. Un normaliseur comporte des bogues, et un champ
dont personne n'avait besoin au moment de son écriture (un indicateur, un nouveau type de message après une mise à niveau de la plateforme) devient nécessaire un an plus tard~; avec la \omterm{def:pl:capturing-and-storing-tick-data:raw}{capture brute}, les deux cas se corrigent
en relançant le normaliseur corrigé~; sans elle, elles sont perdues. Le prix se mesure en octets, que le reste de ce chapitre
compte.

\begin{example}[Un champ que personne n'a conservé]\label{ex:pl:capturing-and-storing-tick-data:field}
Les messages d'exécution avec prix du simulateur (type C) portent un indicateur imprimable qui précise si l'exécution compte dans le
volume de la journée, et ses ajouts d'ordres portent le symbole de l'instrument et un numéro de suivi. L'enregistrement normalisé ne prévoit de champ pour
aucun de ces éléments. Un chercheur qui veut plus tard un volume excluant les exécutions non imprimables ne peut le recalculer
qu'à partir de la \omterm{def:pl:capturing-and-storing-tick-data:raw}{capture brute}~; à partir du stockage normalisé, l'information est perdue pour chaque journée déjà écrite.
\end{example}

\section{Dispositions des enregistrements et codages}

La \Cref{fig:pl:capturing-and-storing-tick-data:layouts} compte les octets par message. La \omterm{def:pl:capturing-and-storing-tick-data:raw}{capture brute} coûte 61.0~octets par message sur une
ligne~: le simulateur envoie un message par paquet, si bien que chaque message porte un en-tête d'enregistrement de fichier de 12~octets, un en-tête MoldUDP64
de 20~octets, une longueur de 2~octets, en plus de ses propres 23 à 44~octets. Les flux réels regroupent plusieurs messages dans un paquet et paient moins de surcharge~;
la part de celle-ci est une propriété du flux, et c'est une raison de mesurer sur le flux effectivement capturé. L'enregistrement normalisé fait
56~octets, de taille fixe~: il est à peine plus petit, mais chaque enregistrement se trouve à un décalage connu, ce qu'exploite le chapitre~4.

\begin{definition}[Taux de compression]\label{def:pl:capturing-and-storing-tick-data:ratio}
Le \emph{taux de compression}\index{taux de compression} (compression ratio) d'une disposition est la taille des données non compressées divisée par la taille des
données compressées, $\mathrm{CR} = b_{\mathrm{raw}} / b_{\mathrm{comp}}$~; les octets par message après compression sont
$b_{\mathrm{raw}} / \mathrm{CR}$.
\end{definition}

\begin{omfigure}
\begin{tikzpicture}
\begin{axis}[width=12cm,height=5.4cm,xbar,bar width=9pt,xmin=0,xmax=72,
  ytick={0,1,2,3,4,5},yticklabels={capture brute,brute~; zlib 6,normalisée,normalisée~; zlib 6,Parquet~; zstd 3,colonnes delta~; zlib 6},
  xlabel={octets par message},tick label style={font=\small},label style={font=\small},y dir=reverse,
  nodes near coords,nodes near coords style={font=\footnotesize,/pgf/number format/fixed,/pgf/number format/precision=1},
  grid=major,grid style={gray!20},table/col sep=comma]
\addplot[fill=omDef!45,draw=omDef] table[x=bytes_per_msg,y=k]{figdata/platforms/02-capturing-and-storing-tick-data/layouts.csv};
\end{axis}
\end{tikzpicture}

\omcaption{Octets par message de la session de dix minutes, par disposition~: la \omterm{def:pl:capturing-and-storing-tick-data:raw}{capture brute} d'une ligne, les enregistrements normalisés de 56~octets,
et chacune compressée par un compresseur générique (zlib, niveau 6), dans un fichier en colonnes avec zstd (chapitre~3), et en colonnes avec les
colonnes de temps, de numéro de séquence et de prix en encodage delta avant zlib. Données~: \texttt{pl\_tickcap.study}.}\label{fig:pl:capturing-and-storing-tick-data:layouts}
\end{omfigure}

Un compresseur générique appliqué à la \omterm{def:pl:capturing-and-storing-tick-data:raw}{capture brute} atteint un taux de 3.32~; sur les enregistrements normalisés, 3.67. Les dispositions en colonnes font mieux,
et le meilleur résultat de ce chapitre est l'astuce en colonnes la plus simple~: écrire chaque champ dans sa propre colonne, et écrire les colonnes dont
valeurs croissent lentement (heure de réception, numéro de séquence, heure d'échange) ou varient par petits pas (prix) sous forme de différences avec la
valeur précédente, avant un compresseur générique. Les deltas du numéro de séquence valent tous 1, ceux de l'heure de réception sont de petits
entiers, et le taux monte à 4.70, soit 11.9~octets par message. Le chapitre~3 explique pourquoi les colonnes se compressent mieux que les lignes~; ce
chapitre se contente de compter.

\omcode{code/firm/tickcap/firm_tickcap.py}{121}{131}{Encodage delta par colonne~: chaque champ écrit comme une colonne, ceux qui varient
lentement sous forme de différences avec la ligne précédente.}\label{lst:pl:capturing-and-storing-tick-data:delta}

\begin{remark}[Répartition des messages]\label{rem:pl:capturing-and-storing-tick-data:mix}
Sur les 258\,535 messages de la session, 49.3~\% sont des ajouts d'ordres et 47.3~\% des suppressions~; les exécutions représentent 3.4~\%. C'est la forme d'un
marché coté, et c'est pourquoi une base ne contenant que les transactions est cent fois plus petite qu'une base à profondeur complète et ne répond à aucune des
questions des chapitres~8 et~18 du livre~7.
\end{remark}

\section{Partitionnement et compression}

\begin{definition}[Partitionnement, clé de partition]\label{def:pl:capturing-and-storing-tick-data:partition}
Le \emph{partitionnement}\index{partitionnement} (partitioning) découpe un jeu de données en parties stockées séparément selon la valeur d'un ou de plusieurs champs, sa
\emph{clé de partition}\index{clé de partition} (partition key), soit la date de négociation ou l'instrument, de sorte qu'une requête qui nomme des valeurs de clé ne lise que
les parties qui les contiennent.
\end{definition}

\begin{proposition}[Élagage]\label{prop:pl:capturing-and-storing-tick-data:pruning}
Une requête dont le filtre fixe la \omterm{def:pl:capturing-and-storing-tick-data:partition}{clé de partition} à un ensemble $S$ de valeurs lit $\sum_{v\in S} B_v$ octets, où $B_v$ est la taille de
la partition de valeur $v$, au lieu de l'ensemble du jeu de données~; avec une surcharge par fichier $o$, elle paie aussi $|S|\,o$.
\end{proposition}

\begin{proof}
Les partitions hors de $S$ ne peuvent contenir aucune ligne correspondante, si bien qu'il n'est pas nécessaire de les ouvrir~; chaque partition ouverte coûte sa taille et sa
surcharge.
\end{proof}

La session du chapitre est stockée sous \texttt{date=2026-09-28/locate=N/}. Ses deux instruments sont très inégaux~: l'instrument~1
occupe 811\,160 octets sur les 14\,477\,960, soit 5.6~\% de la journée, si bien qu'une requête le concernant lit 5.6~\% de ce que coûterait le
fichier non partitionné. Le terme de surcharge est la limite. Partitionner une journée d'options américaines par série d'options, et une requête pour un sous-jacent donné ouvre
des milliers de petits fichiers, chacun coûtant une recherche de répertoire et une ouverture~; partitionner par sous-jacent donne des fichiers volumineux, et la
requête en ouvre peu. La règle empirique qui en découle --- choisir la \omterm{def:pl:capturing-and-storing-tick-data:partition}{clé de partition} parmi les filtres que les requêtes utilisent réellement, et garder
des partitions assez grandes pour que leur surcharge soit négligeable --- fait l'objet de la base de ticks du chapitre~4.

La compression implique son propre compromis, entre taux et vitesse (voir la \cref{fig:pl:capturing-and-storing-tick-data:speed}). Sur cet ordinateur portable
zlib au niveau 9 gagne environ 2~\% de taux par rapport au niveau 6, pour un douzième de sa vitesse, et zstd au niveau 9 dans un fichier en colonnes est
à la fois plus rapide et plus efficace que zlib au niveau 6 sur les mêmes enregistrements. La capture écrit une seule fois et en temps réel, elle a donc besoin de vitesse~; la base historique écrit une fois
la nuit et est lue des milliers de fois, elle privilégie donc le taux.

\begin{omfigure}
\begin{tikzpicture}
\begin{axis}[width=11cm,height=6.8cm,xmode=log,xlabel={débit de compression (Mo/s, échelle logarithmique)},ylabel={taux de compression},
  tick label style={font=\small},label style={font=\small},grid=major,grid style={gray!20},xmin=1.5,xmax=1000,ymin=2.7,ymax=5.0,
  nodes near coords,point meta=explicit symbolic,visualization depends on={value \thisrow{anchor}\as\myanchor},
  nodes near coords style={font=\scriptsize,anchor=\myanchor},table/col sep=comma]
\addplot[only marks,mark=*,omDef] table[x=mb_per_s,y=ratio,meta=method,row sep=\\]{method,ratio,mb_per_s,anchor\\ brut zlib 1,2.949,122.5,south west\\ brut zlib 6,3.318,34.5,north west\\ brut zlib 9,3.399,3.0,south west\\ norm. zlib 6,3.674,51.3,south east\\ Parquet zstd 1,3.454,220.3,south west\\ Parquet zstd 9,3.826,94.2,south east\\ delta zlib 6,4.703,31.2,south west\\};
\end{axis}
\end{tikzpicture}

\omcaption{\omterm{def:pl:capturing-and-storing-tick-data:ratio}{Taux de compression} en fonction du débit pour la \omterm{def:pl:capturing-and-storing-tick-data:raw}{capture brute} de la session («~brute~») et les enregistrements normalisés, par méthode et
niveau, meilleure de cinq exécutions sur un seul cœur. Mesuré sur un ordinateur portable (Intel Core Ultra 7 155H) sous WSL2, machine par ailleurs inactive. Données~:
\texttt{bench\_compress.py}.}\label{fig:pl:capturing-and-storing-tick-data:speed}
\end{omfigure}

\begin{method}[Choix d'une disposition pour une capture]\label{met:pl:capturing-and-storing-tick-data:choose}
Conserver la \omterm{def:pl:capturing-and-storing-tick-data:raw}{capture brute} des lignes sur lesquelles la firme négocie, compressée rapidement, aussi longtemps que les enquêtes en ont besoin. Normaliser en
enregistrements de largeur fixe dans l'ordre d'arrivée pour la base intrajournalière. En fin de journée, réécrire en colonnes, partitionnées par
date et par la clé sur laquelle filtrent les requêtes, compressées pour le \omterm{def:pl:capturing-and-storing-tick-data:ratio}{taux de compression}. Mesurer les octets par message et la vitesse du compresseur sur les
flux propres à la firme, et non sur ceux d'un banc d'essai.
\end{method}

\section{Rétention, niveaux de stockage et économie des pétaoctets}\label{sec:pl:capturing-and-storing-tick-data:economics}

\begin{definition}[Niveau de stockage, politique de rétention]\label{def:pl:capturing-and-storing-tick-data:tier}
Un \emph{niveau de stockage}\index{niveau de stockage} (storage tier) est une catégorie de stockage dotée de son propre prix, de son temps d'accès et de ses frais d'accès (stockage à chaud
lecture en millisecondes, stockage à accès peu fréquent, stockage d'archive restauré en quelques heures). Une \emph{politique de rétention}\index{politique de rétention} (retention policy)
précise pendant combien de temps chaque jeu de données est conservé, dans quel niveau à chaque âge, et quand il est supprimé.
\end{definition}

\begin{dated}{2026-09}{Nombre de messages d'un flux complet et prix du stockage}\label{dat:pl:capturing-and-storing-tick-data:numbers}
L'opérateur du flux consolidé d'options américaines prévoit, dans un avis aux abonnés du 15 septembre 2025 (consulté en
septembre 2026), une capacité de 311 milliards de messages par jour à partir de juillet 2026, avec des pointes de 13.6 millions de messages par 100 millisecondes, pour l'un des
deux flux redondants. Les tarifs publics d'un grand fournisseur de cloud pour le stockage d'objets dans sa région US East (consultés en
septembre 2026) sont de 0.023~USD par Go-mois pour le stockage standard, 0.0125~USD pour l'accès peu fréquent et 0.00099~USD pour son niveau d'archivage le plus profond.
\end{dated}

\begin{dated}{2026-09}{Durée de conservation des enregistrements d'ordres}\label{dat:pl:capturing-and-storing-tick-data:retention}
Dans l'Union européenne (consulté en septembre 2026), les entreprises d'investissement qui pratiquent le trading algorithmique conservent les enregistrements de leurs
ordres pendant cinq ans à compter de la transmission de chaque ordre (Commission Delegated Regulation (EU) 2017/589, article 28). Les données de marché
ne sont pas un enregistrement d'ordre, mais les données dont une firme a besoin pour expliquer un ordre sont conservées au moins aussi longtemps.
\end{dated}

Le modèle de stockage de \texttt{firm.tickcap} le développe~: les octets par jour sont égaux au nombre de messages, multiplié par les octets par message, multiplié par le nombre
de copies~; une année compte 252 jours de séance~; le coût en régime permanent pour conserver cinq ans est celui du dernier mois (21 jours de séance) sur le stockage à chaud
du reste de la première année sur le stockage à accès peu fréquent, et des quatre années plus anciennes sur l'archive.
Le \Cref{tab:pl:capturing-and-storing-tick-data:economics} applique ce modèle à la capacité prévue du flux d'options, avec les
octets par message de la session simulée.

\begin{table}[ht]
\centering\small
\begin{tabular}{lrrrr}
\toprule
disposition & To par jour & Po par an & Po en 5 ans & USD par an, par niveaux \\
\midrule
\omterm{def:pl:capturing-and-storing-tick-data:raw}{capture brute}, deux lignes & 37.95 & 9.56 & 47.8 & 1\,989\,000 \\
\omterm{def:pl:capturing-and-storing-tick-data:raw}{capture brute}, une ligne & 18.97 & 4.78 & 23.9 & 995\,000 \\
normalisée & 17.42 & 4.39 & 21.9 & 913\,000 \\
colonnes delta, zlib 6 & 3.70 & 0.93 & 4.67 & 194\,000 \\
\bottomrule
\end{tabular}
\caption{Capture complète du flux d'options à sa capacité quotidienne prévue (encadré daté), avec les octets par
message de la session simulée~: stockage par jour, par an et après cinq ans, et coût annuel de la conservation de cinq ans sur trois niveaux aux tarifs
datés. Données~: \texttt{pl\_tickcap.economics}.}\label{tab:pl:capturing-and-storing-tick-data:economics}
\end{table}

\begin{omfigure}
\begin{tikzpicture}
\begin{axis}[width=11cm,height=5.6cm,xlabel={années conservées},ylabel={pétaoctets stockés},tick label style={font=\small},
  label style={font=\small},grid=major,grid style={gray!20},xmin=0,xmax=5,ymin=0,ymax=50,xtick={0,1,2,3,4,5},
  legend columns=2,legend style={font=\footnotesize,at={(0.5,-0.28)},anchor=north,draw=none,/tikz/every even column/.append style={column sep=8pt}},
  table/col sep=comma]
\addplot[thick,omThm] table[x=years,y=s0]{figdata/platforms/02-capturing-and-storing-tick-data/stored.csv};
\addplot[thick,omThm,dashed] table[x=years,y=s1]{figdata/platforms/02-capturing-and-storing-tick-data/stored.csv};
\addplot[thick,omDef] table[x=years,y=s2]{figdata/platforms/02-capturing-and-storing-tick-data/stored.csv};
\addplot[thick,omMeth] table[x=years,y=s3]{figdata/platforms/02-capturing-and-storing-tick-data/stored.csv};
\legend{brut~; deux lignes,brut~; une ligne,normalisé,colonnes delta + zlib}
\end{axis}
\end{tikzpicture}

\omcaption{Pétaoctets accumulés par la capture complète du flux d'options à sa capacité prévue, par disposition~: cinq ans de \omterm{def:pl:capturing-and-storing-tick-data:raw}{capture brute} des deux lignes
atteignent 48~Po, cinq ans de colonnes compressées 4.7~Po. Données~: \texttt{pl\_tickcap.economics} (nombre de messages et tarifs de l'encadré daté).
et tarifs de l'encadré daté).}\label{fig:pl:capturing-and-storing-tick-data:stored}
\end{omfigure}

Deux précautions gardent le tableau honnête. Le nombre de messages est une capacité que l'opérateur du flux prévoit, un plafond et non une
journée moyenne, si bien que le tableau constitue une borne supérieure. Et les octets par message sont ceux du simulateur~: les vrais messages d'options ont d'autres
tailles et une autre redondance, et un taux réel peut être supérieur ou inférieur. La méthode se transpose~; les chiffres doivent être mesurés à nouveau sur le
flux lui-même. Ce qui se transpose, c'est la forme~: le pétaoctet le moins cher est celui qui n'est pas écrit deux fois (les deux lignes), pas écrit dans
sa disposition la plus lourde, et pas conservé sur le niveau à chaud après le mois pendant lequel quelqu'un le lit.

\section{Tutoriel~: dix minutes sur deux lignes}\label{tut:pl:capturing-and-storing-tick-data:tut}

\begin{tutorial}
\textbf{Objectif.} Capturer une session simulée chargée sur les deux lignes, arbitrer et normaliser, mesurer chaque disposition, et étendre le
résultat à un flux complet. \textbf{État final~:} \cref{fig:pl:capturing-and-storing-tick-data:layouts,fig:pl:capturing-and-storing-tick-data:stored}
et le \cref{tab:pl:capturing-and-storing-tick-data:economics}.
\begin{tutsteps}
\item \textbf{Générer.} \texttt{pl\_tickcap.day()} lance \texttt{make\_recorded\_day} du livre~10 pendant dix minutes et deux
instruments dans \texttt{data/platforms/generated/} (environ 43~Mo, ignoré par git, régénéré par le script).
\item \textbf{Arbitrer et normaliser} (\cref{lst:pl:capturing-and-storing-tick-data:arb})~: 258\,535 messages, 255\,899
doublons, quatre trous de huit messages.
\item \textbf{Mesurer} avec \texttt{study()}~: octets par message et \omterm{def:pl:capturing-and-storing-tick-data:ratio}{taux de compression} de chaque disposition.
\item \textbf{Partitionner} avec \texttt{firm\_tickcap.partition} et comparer les octets lus par une requête sur un seul instrument.
\item \textbf{Passer à l'échelle} avec \texttt{economics()} et les données sources citées.
\end{tutsteps}
\textbf{Pour aller plus loin.} Trier les enregistrements par instrument avant l'encodage delta, prévoir, puis mesurer, la variation
du taux~; exécuter \texttt{bench\_compress.py} avec zstd à davantage de niveaux et trouver le niveau auquel la capture suivrait le
débit de pointe du flux.
\end{tutorial}

\section{Construction~: la capture des ticks}\label{bld:pl:capturing-and-storing-tick-data:tickcap}

\begin{build}
\bfield{Rôle} Le premier étage de la colonne vertébrale de données de la firme~: \omterm{def:pl:capturing-and-storing-tick-data:raw}{captures brutes} de chaque ligne, enregistrements normalisés et arbitrés, une
disposition partitionnée, et les chiffres qui dimensionnent le stockage. La base de ticks du chapitre~4 lit ses enregistrements~; les contrôles de qualité du chapitre~7
lisent ses trous et ses compteurs.
\bfield{Interface} \texttt{RECORD} (56~octets~: heure de réception et événement de flux du livre~13)~; \texttt{read\_capture(data)}~;
\texttt{arbitrate(lines, timeout\_ns) -> (messages, gaps, counters)}~; \texttt{normalise(messages)}~;
\texttt{partition(rec, root, date)}, \texttt{read\_partition}~; \texttt{delta\_encode},
\texttt{compressed\_size(data, method, level)}~; \texttt{StorageModel(msgs\_per\_day, bytes\_per\_msg, ratio, days\_per\_year, copies)} avec \texttt{bytes\_per\_day},
\texttt{bytes\_per\_year}, \texttt{stored\_after}, \texttt{cost\_per\_year}.
\bfield{Règles} La \omterm{def:pl:capturing-and-storing-tick-data:raw}{capture brute} n'est jamais modifiée~; l'arbitrage est déterministe (première copie par heure de réception, puis par ligne)~; chaque
trou est consigné avec son intervalle de numéros de séquence~; l'enregistrement normalisé est l'événement du livre~13 plus l'heure de réception, octet pour octet.
\bfield{Tests d'acceptation} \texttt{code/firm/tickcap/tests/}~: la disposition de l'enregistrement~; un aller-retour de capture~; l'arbitrage avec une perte sur
une ligne, des doublons et un trou sur les deux~; les mêmes événements que le feed handler du livre~13 sur une session enregistrée jusqu'au premier trou~;
des partitions écrites puis relues~; chaque méthode de compression plus petite que l'entrée et l'encodage delta inversible~; le
modèle de stockage calculé à la main.
\bfield{Pour aller plus loin} La retransmission et la reprise par instantané (snapshot) pendant la capture (réutiliser \texttt{firm.feedhandler})~; une capture qui écrit
des blocs compressés avec un index~; le modèle avec un nombre de messages qui croît chaque année.
\end{build}

\begin{omsources}
\item Options Price Reporting Authority, \emph{Revised OPRA Capacity Projections}, avis aux abonnés aux données multicast,
15 septembre 2025.
\item Y.~Collet et M.~Kucherawy (éd.), \emph{Zstandard Compression and the application/zstd Media Type}, RFC~8878, 2021~;
P.~Deutsch et J.-L.~Gailly, \emph{ZLIB Compressed Data Format Specification}, RFC~1950, 1996.
\item Commission Delegated Regulation (EU) 2017/589 (RTS~6), article~28.
\end{omsources}

\section{Exercices}

\begin{exercise}[$\star$]\label{exo:pl:capturing-and-storing-tick-data:1}
La capture de la ligne A fait 15\,773\,858 octets pour 258\,535 messages, un message par paquet. Combien d'octets par message cela représente-t-il, et
quelle part en est l'encadrement (l'en-tête d'enregistrement de fichier de 12 octets, l'en-tête de paquet de 20 octets et la longueur de 2 octets)~?
\end{exercise}

\begin{exercise}[$\star$]\label{exo:pl:capturing-and-storing-tick-data:2}
Les enregistrements normalisés se compressent avec zlib au niveau 6 à 15.24 octets par message. Quel est le \omterm{def:pl:capturing-and-storing-tick-data:ratio}{taux de compression}~?
\end{exercise}

\begin{exercise}[$\star$]\label{exo:pl:capturing-and-storing-tick-data:3}
La partition de l'instrument~1 contient 811\,160 octets des 14\,477\,960 octets de la journée. Combien d'enregistrements cela représente-t-il, et quelle part de la journée
une requête sur l'instrument~1 lit-elle avec et sans \omterm{def:pl:capturing-and-storing-tick-data:partition}{partitionnement} par instrument~?
\end{exercise}

\begin{exercise}[$\star\star$]\label{exo:pl:capturing-and-storing-tick-data:4}
Conserver les \omterm{def:pl:capturing-and-storing-tick-data:raw}{captures brutes} des deux lignes double leur coût. Formuler une question à laquelle seules les deux \omterm{def:pl:capturing-and-storing-tick-data:raw}{captures brutes} peuvent répondre, et une question
à laquelle les enregistrements arbitrés répondent aussi bien.
\end{exercise}

\begin{exercise}[$\star\star$]\label{exo:pl:capturing-and-storing-tick-data:5}
À raison de 311 milliards de messages par jour et de 11.9 octets par message, combien de téraoctets par jour et de pétaoctets par an (252 jours)
occupe la disposition en colonnes compressée~?
\end{exercise}

\begin{exercise}[$\star\star$]\label{exo:pl:capturing-and-storing-tick-data:6}
Avec les prix datés, combien coûte par an la conservation de cinq années de la disposition compressée si tout reste en
stockage standard, comparé à la politique par niveaux du \cref{tab:pl:capturing-and-storing-tick-data:economics}~?
\end{exercise}

\begin{exercise}[$\star\star\star$]\label{exo:pl:capturing-and-storing-tick-data:7}
\emph{Programmation.} Trier les enregistrements de la séance par instrument puis par numéro de séquence avant l'encodage delta. Prédire la variation
du \omterm{def:pl:capturing-and-storing-tick-data:ratio}{taux de compression}, puis la mesurer avec \texttt{compressed\_size}, et expliquer le résultat.
\end{exercise}

\begin{exercise}[$\star\star\star$]\label{exo:pl:capturing-and-storing-tick-data:8}
\emph{Trouver la faille.} «~Notre capture compresse chaque paquet avec zlib au niveau 9 avant de l'écrire, car cela donne le meilleur
taux sur notre référence.~»
\end{exercise}

\section{Problème~: un pétaoctet par an~?}

\begin{problem}[{Problème du week-end --- dimensionnement de la capture d'un flux complet}]\label{pb:pl:capturing-and-storing-tick-data:1}
La séance de dix minutes du chapitre, ses dispositions, la capacité prévue du flux d'options et les prix de stockage des
encadrés datés.

\medskip\noindent\textbf{Partie I --- La capture.}
\begin{enumerate}
\item Combien de paquets transportent les deux lignes, et combien de messages l'arbitrage délivre-t-il~?
\item Combien de doublons sont écartés, et pourquoi ce nombre est-il proche de la moitié des paquets~?
\item Combien de messages ont été perdus sur les deux lignes, et qu'aurait fait un feed handler en temps réel à leur sujet~?
\item Pourquoi les enregistrements normalisés ne peuvent-ils pas être reconvertis en \omterm{def:pl:capturing-and-storing-tick-data:raw}{capture brute}~?
\item Nommer un champ des messages du simulateur que l'enregistrement normalisé omet.
\end{enumerate}

\medskip\noindent\textbf{Partie II --- Dispositions.}
\begin{enumerate}[resume]
\item Donner les octets par message pour la \omterm{def:pl:capturing-and-storing-tick-data:raw}{capture brute}, pour les enregistrements normalisés et pour les colonnes soumises à l'encodage delta avec zlib.
\item Quels sont les \omterm{def:pl:capturing-and-storing-tick-data:ratio}{taux de compression} de la \omterm{def:pl:capturing-and-storing-tick-data:raw}{capture brute} avec zlib et des colonnes soumises à l'encodage delta~?
\item Pourquoi l'encodage delta profite-t-il le plus au numéro de séquence~?
\item Quelle disposition faudrait-il conserver, et pendant combien de temps~?
\item Pourquoi faut-il remesurer les octets par message sur le flux réel~?
\end{enumerate}

\medskip\noindent\textbf{Partie III --- Échelle.}
\begin{enumerate}[resume]
\item À la capacité prévue, combien de téraoctets par jour représente la \omterm{def:pl:capturing-and-storing-tick-data:raw}{capture brute} des deux lignes~?
\item Combien de pétaoctets s'accumulent-ils en cinq ans~?
\item Mêmes questions pour les colonnes soumises à l'encodage delta.
\item Quel est le coût annuel de la conservation de cinq années de chacune dans les trois niveaux~?
\item Pourquoi le chiffre de capacité est-il une borne supérieure pour une journée typique~?
\end{enumerate}

\medskip\noindent\textbf{Partie IV --- Le verdict.}
\begin{enumerate}[resume]
\item Énoncer le \emph{résultat nommé}~: les pétaoctets par an et le coût annuel de la capture complète dans la disposition la plus lourde et dans la plus
légère.
\item La réponse à la question du titre dépend-elle davantage de la disposition ou de la hiérarchisation par niveaux~?
\item Que fait gagner le \omterm{def:pl:capturing-and-storing-tick-data:partition}{partitionnement} par instrument à une requête sur un instrument pendant la séance, et que coûte-t-il
lorsqu'il existe des millions de séries~?
\item À quel endroit du pipeline faut-il exécuter le compresseur rapide, et à quel endroit le lent~?
\item En une phrase~: que doit conserver une firme de ses \omterm{def:pl:capturing-and-storing-tick-data:tick}{données tick}, et sous quelle forme~?
\end{enumerate}
\end{problem}

\section{Questions d'entretien}

\begin{interviewq}[$\star$ \iqroles{développeur, chercheur}]\label{iq:pl:capturing-and-storing-tick-data:1}
Quelle est la différence entre les \omterm{def:pl:capturing-and-storing-tick-data:tick}{données tick} brutes et normalisées, et pourquoi une firme conserverait-elle les deux~?
\end{interviewq}

\begin{interviewq}[$\star$ \iqroles{développeur}]\label{iq:pl:capturing-and-storing-tick-data:2}
Deux lignes de flux redondantes sont capturées. Comment les transformez-vous en un seul flux de messages, et que faites-vous d'un trou~?
\end{interviewq}

\begin{interviewq}[$\star\star$ \iqroles{développeur}]\label{iq:pl:capturing-and-storing-tick-data:3}
Comment estimeriez-vous le stockage nécessaire pour conserver cinq années d'un flux à profondeur complète, et que mesureriez-vous en premier~?
\end{interviewq}

\begin{interviewq}[$\star\star$ \iqroles{développeur}]\label{iq:pl:capturing-and-storing-tick-data:4}
Comment choisissez-vous une \omterm{def:pl:capturing-and-storing-tick-data:partition}{clé de partition} pour les \omterm{def:pl:capturing-and-storing-tick-data:tick}{données tick}, et que se passe-t-il si vous en choisissez une trop fine~?
\end{interviewq}

\begin{interviewq}[$\star\star$ \iqroles{développeur, chercheur}]\label{iq:pl:capturing-and-storing-tick-data:5}
Pourquoi les données de marché se compressent-elles bien, et quels champs se compressent le mieux~?
\end{interviewq}

\begin{interviewq}[$\star\star\star$ \iqroles{développeur}]\label{iq:pl:capturing-and-storing-tick-data:6}
Concevez le stockage des \omterm{def:pl:capturing-and-storing-tick-data:tick}{données tick} d'une firme qui négocie des options américaines~: ce qui est capturé où, dans quelle disposition, conservé combien de temps, dans
quel niveau, et pour quel ordre de grandeur de coût~?
\end{interviewq}
